\documentclass[12pt]{article}
\usepackage{cocina1}

\begin{document}
\begin{ficha}{Comparar}{¿Quién se lleva más?}{2}

\begin{encargo}
Dos mesas, dos pedidos. \textbf{Apuesta} primero en la casilla ($>$, $<$ o $=$) y luego compruébalo.
\end{encargo}

\bloque{1}{Mira cómo se hace}{Con pizza}

\vspace{1.5mm}
\begin{minipage}[c]{.55\linewidth}
\paso{1}\textbf{¿Los trozos son iguales?}\ Sí: octavos los dos.\\[1.5mm]
\paso{2}\textbf{Cuenta los trozos.}\ 3 y 5.\\[1.5mm]
\paso{3}\textbf{Pon el signo.}\quad \fA{3}{8}\ $<$\ \fB{5}{8}
\end{minipage}%
\begin{minipage}[c]{.45\linewidth}\centering
\pizza[escala=.85]{8}{0,1,2}\qquad\pizza[escala=.85]{8}{0,1,2,3,4}
\end{minipage}

\vspace{1.5mm}
\begin{nick}
Si los trozos son \textbf{iguales}, cuenta. Si \textbf{no} lo son, primero
\textbf{córtalos igual}, y luego cuenta.
\end{nick}

\bloque{2}{Ahora tú}{Con pizza}\quad{\small colorea los pedidos y pon el signo}

\vspace{1mm}
\begin{minipage}[t]{.33\linewidth}\vspace{0pt}
\textbf{a)} Trozos iguales\\[1mm]
\pizza[escala=.78]{6}{}\ \pizza[escala=.78]{6}{}\\[1mm]
\fA{2}{6}\ \signo\ \fB{5}{6}
\end{minipage}%
\begin{minipage}[t]{.33\linewidth}\vspace{0pt}
\textbf{b)} Un trozo cada uno\\[1mm]
\pizza[escala=.78]{2}{}\ \pizza[escala=.78]{4}{}\\[1mm]
\fA{1}{2}\ \signo\ \fB{1}{4}
\end{minipage}%
\begin{minipage}[t]{.34\linewidth}\vspace{0pt}
\textbf{c)} Corta la de cuartos\\[1mm]
\pizza[escala=.78,guias=8]{4}{}\ \pizza[escala=.78]{8}{}\\[1mm]
\fA{3}{4}\ \signo\ \fB{5}{8}\quad $\dfrac{3}{4} = \fhueco[][8]$
\end{minipage}

\alto{el 4 de $\frac{1}{4}$ es más grande que el 2 de $\frac{1}{2}$, pero el
cuarto es \textbf{más pequeño}: el número de abajo dice en cuántos trozos se
corta.}

\reto{Nick y Migas tienen dos pizzas iguales. Nick se come $\frac{2}{4}$ y
Migas $\frac{3}{6}$. ¿Quién ha comido más?}

\comprueba{Cada código abre esas dos mesas.}{%
\qr{1}{j=comparar\&a=3/8\&b=5/8}\hfill\qr{2a}{j=comparar\&a=2/6\&b=5/6}\hfill
\qr{2b}{j=comparar\&a=1/2\&b=1/4}\hfill\qr{2c}{j=comparar\&a=3/4\&b=5/8}\hfill
\qr{reto}{j=comparar\&a=2/4\&b=3/6}}

\end{ficha}

% ─────────────────────────────── REVERSO ──────────────────────────────────
\begin{dorso}{Comparar}{Más mesas}{2}

\bloque{1}{En la regla}{Con la regla}\quad{\small marca las dos y pon el signo}

\vspace{2mm}
\textbf{a)}\ $\dfrac{1}{2}$\ \signo\ $\dfrac{3}{8}$\qquad\reglavacia[10cm]{1}{8}

\vspace{3mm}
\textbf{b)}\ $\dfrac{3}{4}$\ \signo\ $\dfrac{5}{8}$\qquad\reglavacia[10cm]{1}{8}

\begin{minipage}[t]{.6\linewidth}\vspace{0pt}
\bloque{2}{Sin pizza}{Sobre el papel}

\vspace{2mm}
{\renewcommand{\arraystretch}{2.5}
\begin{tabular}{@{}l@{\hspace{3mm}}l@{}}
\textbf{a)}\ $\dfrac{2}{5}\ \signo\ \dfrac{4}{5}$ & \\
\textbf{b)}\ $\dfrac{2}{3}\ \signo\ \dfrac{5}{6}$ & $\dfrac{2}{3} = \fhueco[][6]$\\
\textbf{c)}\ $\dfrac{1}{3}\ \signo\ \dfrac{1}{5}$ & \\
\textbf{d)}\ $\dfrac{3}{4}\ \signo\ \dfrac{7}{12}$ & $\dfrac{3}{4} = \fhueco[][12]$\\
\end{tabular}}
\end{minipage}\hfill
\begin{minipage}[t]{.37\linewidth}\vspace{6mm}
\tablamult
\end{minipage}

\bloque{3}{Problemas}{Sobre el papel}

\vspace{1.5mm}
\textbf{a)} En la mesa 4 se comen $\frac{2}{3}$ de pizza y en la mesa 7,
$\frac{4}{6}$ de otra igual. ¿Quién come más?\hfill\hueco[2cm]

\vspace{3mm}
\textbf{b)} Ana corre $\frac{3}{4}$ de kilómetro y Luis $\frac{5}{8}$. ¿Quién
corre más?\hfill\hueco[2cm]

\vspace{3mm}
\begin{semaforo}
\sfila{Comparo si los trozos son iguales}
\sfila{Corto una pizza para que los trozos sean iguales}
\sfila{Pongo bien el signo $<$, $>$ o $=$}
\end{semaforo}

\comprueba{Cada código abre esas dos mesas.}{%
\qr{2a}{j=comparar\&a=2/5\&b=4/5}\hfill\qr{2b}{j=comparar\&a=2/3\&b=5/6}\hfill
\qr{2c}{j=comparar\&a=1/3\&b=1/5}\hfill\qr{2d}{j=comparar\&a=3/4\&b=7/12}\hfill
\qr{3a}{j=comparar\&a=2/3\&b=4/6}\hfill\qr{3b}{j=comparar\&a=3/4\&b=5/8}}
\end{dorso}
\end{document}
